% 统计图真值来自 source/普通高考/2010/2010湖南理.pdf 第 1 页、题 17。
% bars: [0,1) height x=0.12, [1,2) height 0.37, [2,3) height 0.39,
%       [3,4) height 0.10, [4,5) height 0.02；组距均为 1。
% dashed guides: y=0.39 to x=3, y=0.37 to x=2, y=0.10 to x=4, y=0.02 to x=5。
% labels: y-axis is 频率/组距, x-axis is 月均用水量/吨, first bar height is x。
\begin{tikzpicture}
\begin{axis}[exam statistical histogram,scaled ticks=false,
  width=16cm,
  height=13.172cm,
  xmin=0, xmax=6,
  ymin=0, ymax=0.45,
  axis lines=left,
  axis line style={-stealth},
  xtick={0,1,2,3,4,5},
  ytick=\empty,
  tick style={draw=none},
  clip=false,
]
    \examstatxlabel{月均用水量/吨};
    \examstatylabel{\(\dfrac{\text{频率}}{\text{组距}}\)};
\addplot[
  ybar interval,
  draw=black,
  fill=none,
  line width=0.45pt,
] coordinates {
  (0,0.12) (1,0.37) (2,0.39) (3,0.10) (4,0.02) (5,0)
};
\draw[dashed,line width=0.4pt] (axis cs:0,0.39) -- (axis cs:3,0.39);
\draw[dashed,line width=0.4pt] (axis cs:0,0.37) -- (axis cs:2,0.37);
\draw[dashed,line width=0.4pt] (axis cs:0,0.10) -- (axis cs:4,0.10);
\draw[dashed,line width=0.4pt] (axis cs:0,0.02) -- (axis cs:5,0.02);
\node[anchor=east,xshift=-4pt] at (axis cs:0.000000,0.39) {0.39};
\node[anchor=east,xshift=-4pt] at (axis cs:0.000000,0.37) {0.37};
\node[anchor=east,xshift=-4pt] at (axis cs:0.000000,0.10) {0.1};
\node[anchor=east,xshift=-4pt] at (axis cs:0.000000,0.02) {0.02};
\node[anchor=east,xshift=-3pt] at (axis cs:0,0.12) {$x$};
\end{axis}
\end{tikzpicture}
